
	% Glossareinträge
\newglossaryentry{Human-in-the-Loop}{
	name=Human in the Loop,
	description={Der Mensch ist aktiv am Entscheidungsprozess beteiligt, indem er mit der simulierten oder automatisierten Umgebung interagiert und diese beeinflusst}
}

\newglossaryentry{CARLA}{
	name=CARLA,
	description={Open-Source-Simulationsumgebung für autonomes Fahren und Verkehrsszenarien}
}

\newglossaryentry{Instance Segmentation}{
	name=Instance Segmentation,
	description={Verfahren zur Erkennung und Klassifizierung einzelner Objekte in digitalen Bildern}
}

\newglossaryentry{Server-Client Struktur}{
	name=Server-Client Struktur,
	description={Architekturmodell, bei dem ein Server zentrale Ressourcen verwaltet und Clients darauf zugreifen}
}

\newglossaryentry{Bildwiederholrate}{
	name=Bildwiederholrate,
	description={Wichtiges Kriterium bei Kameras oder Bildschirmen. Gibt an wie viele Bilder pro Sekunde aufgenommen bzw. dargestellt werden}
}

\newglossaryentry{Open Source}{
	name=Open Source,
	description={Softwarelizenzmodell – hilfreich zur Abgrenzung gegenüber kommerziellen Systemen}
}

\newglossaryentry{Actor}{
	name=Actor,
	description={Oberbegriff für Sensoren, Fahrzeuge etc. im CARLA-Ökosystem}
}

\newglossaryentry{Lidar}{
	name=Lidar,
	description={Alternativer Sensortyp basierend auf Lasertechnologie}
}

\newglossaryentry{Pixel}{
	name=Pixel,
	description={Kleinste Einheit eines digitalen Bildes, oft in Bezug auf Bildschirmkoordinaten}
}

\newglossaryentry{Raw-Output}{
	name=Raw-Output,
	description={Unverarbeitete Ausgabedaten des Eye Trackers}
}

\newglossaryentry{e-Funktion}{
	name=e-Funktion,
	description={Mathematische Exponentialfunktion, verwendet zur Modellierung von Anstiegsverhalten}
}

\newglossaryentry{quadratische Funktion}{
	name=Quadratische Funktion,
	description={Mathematische Exponentialfunktion, verwendet zur Modellierung von Anstiegsverhalten}
}

\newglossaryentry{V-Funktion}{
	name=V-Funktion	,
	description={Mathematische Exponentialfunktion, verwendet zur Modellierung von Anstiegsverhalten}
}

\printglossaries
